\section{Introduction}
Zero-knowledge proofs demonstrate that a computation is correct without revealing its private inputs. There are two types of zero-knowledge proof system. In a zkSNARK~\cite{groth2016}, the computation to prove is defined before the system is built. Conversely, in a zkVM~\cite{yang2026sokzkvm} the computation is an input, so it is flexible. An anonymous credential lets a holder prove a statement about themselves without revealing who they are. What a holder must prove differs from one verifier to the next, so a zkVM is more suitable. Showing such a credential requires verifying a signature, and this thesis uses a post-quantum signature scheme~\cite{10.1007/978-981-95-2961-2_6}. However, proving a signature's verification inside a zkVM is expensive. One signature verification did not finish in five hours on a personal laptop. This thesis makes two contributions. First, a cost criterion for signature verification. Before building a program, it decides whether a zkVM suffices or a precompile is needed to speed up the costly operations of the major primitives. Second, we run BDEC~\cite{10.1007/978-981-96-0957-4_3}, a post-quantum anonymous credential, inside a zkVM, and measure its cost and which of its security guarantees survive. To our knowledge, this is the first post-quantum anonymous credential run inside a zkVM. The zkVM we study cannot make a single proof that is both post-quantum and zero-knowledge. Achieving both in one proof is left for future work.